% Options for packages loaded elsewhere
\PassOptionsToPackage{unicode}{hyperref}
\PassOptionsToPackage{hyphens}{url}
\documentclass[
]{article}
\usepackage{xcolor}
\usepackage[margin=1in]{geometry}
\usepackage{amsmath,amssymb}
\setcounter{secnumdepth}{-\maxdimen} % remove section numbering
\usepackage{iftex}
\ifPDFTeX
  \usepackage[T1]{fontenc}
  \usepackage[utf8]{inputenc}
  \usepackage{textcomp} % provide euro and other symbols
\else % if luatex or xetex
  \usepackage{unicode-math} % this also loads fontspec
  \defaultfontfeatures{Scale=MatchLowercase}
  \defaultfontfeatures[\rmfamily]{Ligatures=TeX,Scale=1}
\fi
\usepackage{lmodern}
\ifPDFTeX\else
  % xetex/luatex font selection
\fi
% Use upquote if available, for straight quotes in verbatim environments
\IfFileExists{upquote.sty}{\usepackage{upquote}}{}
\IfFileExists{microtype.sty}{% use microtype if available
  \usepackage[]{microtype}
  \UseMicrotypeSet[protrusion]{basicmath} % disable protrusion for tt fonts
}{}
\makeatletter
\@ifundefined{KOMAClassName}{% if non-KOMA class
  \IfFileExists{parskip.sty}{%
    \usepackage{parskip}
  }{% else
    \setlength{\parindent}{0pt}
    \setlength{\parskip}{6pt plus 2pt minus 1pt}}
}{% if KOMA class
  \KOMAoptions{parskip=half}}
\makeatother
\usepackage{longtable,booktabs,array}
\newcounter{none} % for unnumbered tables
\usepackage{calc} % for calculating minipage widths
% Correct order of tables after \paragraph or \subparagraph
\usepackage{etoolbox}
\makeatletter
\patchcmd\longtable{\par}{\if@noskipsec\mbox{}\fi\par}{}{}
\makeatother
% Allow footnotes in longtable head/foot
\IfFileExists{footnotehyper.sty}{\usepackage{footnotehyper}}{\usepackage{footnote}}
\makesavenoteenv{longtable}
\usepackage{graphicx}
\makeatletter
\newsavebox\pandoc@box
\newcommand*\pandocbounded[1]{% scales image to fit in text height/width
  \sbox\pandoc@box{#1}%
  \Gscale@div\@tempa{\textheight}{\dimexpr\ht\pandoc@box+\dp\pandoc@box\relax}%
  \Gscale@div\@tempb{\linewidth}{\wd\pandoc@box}%
  \ifdim\@tempb\p@<\@tempa\p@\let\@tempa\@tempb\fi% select the smaller of both
  \ifdim\@tempa\p@<\p@\scalebox{\@tempa}{\usebox\pandoc@box}%
  \else\usebox{\pandoc@box}%
  \fi%
}
% Set default figure placement to htbp
\def\fps@figure{htbp}
\makeatother
\setlength{\emergencystretch}{3em} % prevent overfull lines
\providecommand{\tightlist}{%
  \setlength{\itemsep}{0pt}\setlength{\parskip}{0pt}}
\usepackage{bookmark}
\IfFileExists{xurl.sty}{\usepackage{xurl}}{} % add URL line breaks if available
\urlstyle{same}
\hypersetup{
  pdftitle={Nuclei Isolation Protocol},
  pdfauthor={Kevin Shih},
  hidelinks,
  pdfcreator={LaTeX via pandoc}}

\title{Nuclei Isolation Protocol}
\author{Kevin Shih}
\date{17 July 2026}

\begin{document}
\maketitle

{
\setcounter{tocdepth}{5}
\tableofcontents
}
\section{Nuclei Isolation Protocol
Materials}\label{nuclei-isolation-protocol-materials}

\begin{itemize}
\tightlist
\item
  10X Genomics Chromium Nuclei Isolation Kit
\item
  Tissue scissors
\item
  Beaker for liquid waste
\item
  0.4\% Trypan Blue
\end{itemize}

\subsection{Lysis Buffer}\label{lysis-buffer}

{\def\LTcaptype{none} % do not increment counter
\begin{longtable}[]{@{}lllll@{}}
\toprule\noalign{}
Reagent & 1 Reaction & 2 Reactions & 3 Reactions & 4 Reactions \\
\midrule\noalign{}
\endhead
\bottomrule\noalign{}
\endlastfoot
Lysis Reagent & 536uL & 1mL 72uL & 1mL 608uL & 2mL 144uL \\
Reducing Agent B & 0.5uL & 1uL & 1.5uL & 2uL \\
Surfactant A & 5.5uL & 11uL & 16.5uL & 22uL \\
RNase Inhibitor & 13.9uL & 27.8uL & 41.7uL & 55.6uL \\
\end{longtable}
}

Note: Lysis Reagent and Surfactant A are found in the 4°C.\\
Note: Reducing Agent B and RNase Inhibitor are found in the -20°C.

\subsection{Debris Removal Buffer}\label{debris-removal-buffer}

{\def\LTcaptype{none} % do not increment counter
\begin{longtable}[]{@{}lllll@{}}
\toprule\noalign{}
Reagent & 1 Reaction & 2 Reactions & 3 Reactions & 4 Reactions \\
\midrule\noalign{}
\endhead
\bottomrule\noalign{}
\endlastfoot
Debris Removal & 550uL & 1mL 100uL & 1mL 650uL & 2mL 200uL \\
Reducing Agent B & 0.5uL & 1uL & 1.5uL & 2uL \\
\end{longtable}
}

Note: Debris Removal Reagent is found in the 4°C.\\
Note: Reducing Agent B is found in the -20°C.

\subsection{Wash and Resuspension
Buffer}\label{wash-and-resuspension-buffer}

{\def\LTcaptype{none} % do not increment counter
\begin{longtable}[]{@{}lllll@{}}
\toprule\noalign{}
Reagent & 1 Reaction & 2 Reactions & 3 Reactions & 4 Reactions \\
\midrule\noalign{}
\endhead
\bottomrule\noalign{}
\endlastfoot
1X PBS & 2mL 887.5uL & 5mL 775uL & 8mL 662.5uL & 11mL 550uL \\
10\% BSA & 330uL & 660uL & 990uL & 1mL 320uL \\
RNase Inhibitor & 82.5uL & 165uL & 247.5uL & 330uL \\
\end{longtable}
}

Note: 1X PBS is found in Bench 3 cabinet.\\
Note: 10\% BSA is found in the 4°C.

\section{Step-by-Step Protocol}\label{step-by-step-protocol}

\subsection{Experimental Setup}\label{experimental-setup}

\begin{enumerate}
\def\labelenumi{\arabic{enumi}.}
\tightlist
\item
  Obtain ice (make sure to keep all samples and reagents on ice)
\item
  Prepare and label the reagents in 15mL conical tubes and keep on ice.
\item
  Label 2 round bottom-tubes from the 10X Consumables Box and 1
  pointed-bottom tubes and keep on ice.
\item
  Obtain the tube with the half-kidney from the -80°C freezer.
\item
  Turn on the centrifuge and set to 4°C for fast temperature control
\end{enumerate}

\subsection{Experiment}\label{experiment}

\begin{enumerate}
\def\labelenumi{\arabic{enumi}.}
\setcounter{enumi}{5}
\tightlist
\item
  Transfer the sample from the tube in the -80°C to the new
  pointed-bottom tube that was obtained from the Consumables Box.
\item
  Add 200uL of the Lysis Buffer to the sample
\item
  Using surgical scissors, mince the 1/2 kidney as much as possible to
  prevent large chunks of tissue being present
\item
  Using the green pestle, \emph{GENTLY} push down against the tissue for
  a maximum of 10x
\item
  Add 300uL of Lysis Buffer and incubate for 10 minutes on ice\\
  10a. During this time, set up your centrifuge to 16,000 rcf (also may
  show up as 16,000 g; rcf=g), for 20 seconds, at 4°C.\\
  10b. Set up the Collection Tube (round-bottom) with the Nuclei
  Isolation Column in it on ice during the lysis so the tube can cool
\item
  Transfer the dissociated tissue onto the column and close the top of
  the Nuclei Isolation Column
\item
  Put the tube into the centrifuge and spin at the preset settings
  (16,000rcf, 20s, 4°C) (make sure to balance centrifuge)
\item
  Once the centrifuge is done, discard the column (pellet should be
  visible at the bottom/side of the tube).
\item
  If tube is fully intact, cap the tube and vortex for 10 seconds to
  resuspend the nuclei. If cap is not intact see 14a\\
  14a. If the cap has broken off during the centrifuge step (step 12),
  resuspend your nuclei by pipetting gently up and down for a maximum of
  10x (you should see a curtain effect where the cells lift off the edge
  of the cap and float initially). Make sure to fully resuspend the
  cells so no clumps are present
\item
  Set your centrifuge to 500 rcf (g), for 3 minutes, at 4°C and spin
  your sample down (make sure to balance centrifuge and put the spine of
  the cap facing outwards)
\item
  Gently remove your sample from the centrifuge and close the centrifuge
  to keep it cold
\item
  Carefully pipette your supernatant (excess fluid) into a waste
  container (your pellet should be visible so it will be easier to see
  and should be on the side of the tube's spine)
\item
  Resuspend the pellet in 500uL of Debris Removal Buffer and place on
  ice
\item
  Change the settings on the centrifuge to 700rcf, for 10 minutes, at
  4°C and spin your sample (balance the centrifuge)
\item
  Carefully remove the sample from the centrifuge (note that your pellet
  will have seemingly disappeared)
\item
  Carefully remove as much of the Debris Removal Buffer as possible
  without disturbing the pellet (if the pellet starts to move then do
  NOT continue removing buffer)
\item
  Resuspend the nuclei with 1mL of Resuspension Buffer (can pipette up
  and down for a maximum of 10x)
\item
  Set the centrifuge to 500rcf, for 5 min, at 4°C and spin down cells
  (balance the centrifuge)
\item
  Remove your sample from the centrifuge carefully (pellet should be
  visible), and remove the supernatant
\item
  Resuspend the nuclei with 1mL of Resuspension Buffer (can pipette up
  and down for a maximum of 10x)
\item
  Spin sample again at 500rcf, for 5 min, at 4°C (balance the
  centrifuge)
\item
  Remove your sample from the centrifuge carefully (pellet should be
  visible), and remove the supernatant
\item
  Resuspend your pellet in 20-500uL of Wash and Resuspension Buffer
\item
  Vortex for 3 seconds
\end{enumerate}

\subsection{Visualization}\label{visualization}

\subsubsection{Protocol}\label{protocol}

\begin{enumerate}
\def\labelenumi{\arabic{enumi}.}
\setcounter{enumi}{29}
\tightlist
\item
  Obtain a 0.6mL Eppendorf Tube
\item
  Add 10uL of 0.4\% Trypan Blue to the bottom of the tube
\item
  Add 10uL of Sample into the bottom of the tube and gently pipette up
  and down a few times to mix
\item
  Incubate on ice for a minimum of 2 minutes
\item
  Pipette 10uL of the dyed sample onto a hemocytometer (with the cover
  slip on top) and allow for capillary action to move the sample evenly
\end{enumerate}

\subsubsection{Light Microscopy}\label{light-microscopy}

Note: Located in the 2nd lab

\begin{enumerate}
\def\labelenumi{\arabic{enumi}.}
\setcounter{enumi}{34}
\tightlist
\item
  Place slide on stage of Leica DMIL LED light microscope
\item
  Turn on the light (right side of the microscope)
\end{enumerate}

\subsubsection{Epifluorescence
Microscopy}\label{epifluorescence-microscopy}

Note: Located in the microscope room

\begin{enumerate}
\def\labelenumi{\arabic{enumi}.}
\setcounter{enumi}{36}
\tightlist
\item
  Log in to the computer
\item
  Open the LASX program, wait for it to load, and click ``Ok''
\item
  To the left of the microscope, turn the switch on for the CiR Advanced
\item
  To the box to the left of the CiR Advanced, turn on the switch
\item
  Turn on the shutter on the main screen and click on the objective with
  the eye to allow light to reach the eye piece
\item
  Once nuclei are in focus, on the computer click the ``Live'' button on
  the lower left part of the screen
\item
  Adjust Gain and Exposure levels in the upper left part of the LASX
  software
\item
  To take the picture, click on ``Single Image'' which will take the
  picture and immediately turn off Live imaging
\item
  Under projects, save the picture/folder 45a. For individual pictures,
  right click and export to either .png, .tif, .jpeg
\item
  To turn off the machine, close the LASX software completely
\item
  Turn off the light box to the left of the CiR Advanced
\item
  Turn off CiR Advanced
\item
  Remove your slide
\item
  Clean station, log out of computer, and close door
\end{enumerate}

This protocol was last amended on July 17th, 2026 by Kevin Shih

\end{document}
